%!TEX root = ../../main.tex
\section{향후 과제}
\label{sec:disc-future}

다음 항목은 모두 이연된 것이며 어느 것도 실행되지 않았고 증거가 아니다. 사전 등록 형식의 계획은 부록 \ref{app:deferred}에 있다.

\begin{itemize}[leftmargin=1.5em]
  \item \textbf{F1 외부 확률 어댑터.} 새 환경의 작은 보류 조각에서 적합한 저차원 단조 어댑터가 $\PTflex$에 적용한 같은 어댑터보다 Brier를 줄이는지, $\Vhat$·$q$·$\PTflex$의 가중치를 동결한 채 검정한다. 경쟁 가설(확률 척도 이동, 평가기 자체의 이동, 구성 이동, 인공물)과 그 구별 예측을 함께 둔다.
  \item \textbf{F2 참여 인원 셀.} $\nmin = 2 / 3 / 4 / 5+$ 셀 안에서 기존 예측표를 재점수화하여 코호트 안 이득의 분포를 기술한다. ``기울기'' 문구는 금지된다.
  \item \textbf{F3 공유 $\PTflexS$ 객체.} S arm 사이에 하나의 동결 $\PTflexS$를 재사용해도 대비가 변하지 않음을 확인하는 위생 항목이다.
  \item \textbf{F4 예비 선언된 T 대 S 외부 전이 대비.} 두 코호트의 외부 $\dBrier$ 차이가 경기 표본 변동보다 큰지를 같은 외부 경기에서 짝 부트스트랩으로 검정한다. 기전 주장은 계속 금지된다.
  \item \textbf{F5 정의 상수 민감도.} $G$, $D$를 부트스트랩 구간 안에서, $B$를 10초와 20초로 옮겨 고정된 2만 경기 부표본에서 재검출·재라벨하고 주 대비를 설정별로 보고한다. 단위의 강건성 검사이지 모델의 검사가 아니다.
  \item \textbf{F6 정보군 모듈 (M-F0--M-F3).} 동일 행·동일 학습기의 중첩 모델 사이 짝 $\dBrier$로 초기 승률 요약에 조건부인 특징군의 추가 효용을 잰다. M-F0는 ``$C$-only 로지스틱''이며 PT 기준선이 아니다.
  \item \textbf{사례 추적.} 소수의 사례를 원자료의 timestamp $\to$ 교전 경계 $\to$ 사용 가능한 $X$ $\to$ 두 시점의 $\Vhat$ $\to$ $\SVI$ $\to$ $q$ 예측 $\to$ 결과 대응으로 연결하는 서사를 본문화한다. 대표성의 증거로 과대 해석하지 않는다.
\end{itemize}

PUBG 등 다른 게임과의 통합, 모든 그래프·시계열 구조의 재실험, 실시간 교전 발생 탐지, 인과적 최적 행동 추천, 새 수작업 주석 연구, 연속형 $\dV$로의 주 목표 교체는 이 학위논문의 범위에 자동으로 포함되지 않는다.
